\chapter{Site map}
\label{cap:sitemap}

A aplicação divide-se numa área pública, com a entrada, e numa área autenticada com uma secção por perfil. Depois do login, cada utilizador é levado para a página inicial do seu perfil e só vê os ecrãs dessa secção, mais os que são comuns a todos (Figura~\ref{fig:sitemap}).

\begin{figure}[H]
    \centering
    \resizebox{\linewidth}{!}{% Site map: uma coluna por área de acesso; cada nó é um ecrã, com a rota por baixo.
% Linha cheia = navegação hierárquica; tracejado = atalho entre ecrãs.
\begin{tikzpicture}[ecra/.append style={text width=2.9cm}]
  % ---------- Área pública ----------
  \node[nomearea=cPublico, anchor=west] (lp) at (-1.55,0.95) {Área pública};
  \node[ecra=cPublico] (p1) at (0,0)    {Entrada e login\rota{/ \ \ /signin}};
  \node[ecra=cPublico] (p2) at (0,-1.5) {Página não encontrada\rota{/not-found}};

  % ---------- Administrador ----------
  \node[nomearea=cAdmin, anchor=west] (la) at (2.35,0.95) {Administrador};
  \node[ecra=cAdmin] (a1) at (3.9,0)    {Painel de administração\rota{/admin/dashboard}};
  \node[ecra=cAdmin] (a2) at (3.9,-1.5) {Gestão de utilizadores\rota{/admin/users}};
  \node[ecra=cAdmin] (a3) at (3.9,-3.0) {Moldes de percurso\rota{/templates}};
  \node[ecra=cAdmin] (a4) at (3.9,-4.5) {Construtor, modo Molde\rota{(dentro de /templates)}};

  % ---------- Supervisor ----------
  \node[nomearea=cSupervisor, anchor=west] (ls) at (6.25,0.95) {Supervisor};
  \node[ecra=cSupervisor] (s1) at (7.8,0)    {Painel do supervisor\rota{/dashboard}};
  \node[ecra=cSupervisor] (s2) at (7.8,-1.5) {Gestão do percurso e inscritos\rota{/pathway/\{id\}/students}};
  \node[ecra=cSupervisor] (s3) at (7.8,-3.0) {Ficha de avaliação\rota{/evaluate/\{id\}}};
  \node[ecra=cSupervisor] (s4) at (7.8,-4.5) {Construtor, modos Geral e Individual\rota{/supervisor/}\rota{pathway-builder/\ldots{}}};
  \node[ecra=cSupervisor] (s5) at (7.8,-6.0) {Moldes de percurso\rota{/templates}};

  % ---------- Supervisionado ----------
  \node[nomearea=cAluno, anchor=west] (lu) at (10.15,0.95) {Supervisionado};
  \node[ecra=cAluno] (u1) at (11.7,0)    {Os meus percursos\rota{/my-profile}};
  \node[ecra=cAluno] (u2) at (11.7,-1.5) {Detalhe do percurso\rota{/pathway/\{id\}}};
  \node[ecra=cAluno] (u3) at (11.7,-3.0) {Detalhe do módulo e documentos\rota{/module/\{id\}}};
  \node[ecra=cAluno] (u4) at (11.7,-4.5) {Conquistas\rota{/badges}};
  \node[ecra=cAluno, dashed] (u5) at (11.7,-6.0) {Relatório de avaliação (PDF)\rota{/api/evaluation/\ldots{}/pdf}};

  % ---------- Comum aos perfis autenticados ----------
  \node[nomearea=cComum, anchor=west] (lc) at (2.35,-7.45) {Comum aos três perfis (cabeçalho)};
  \node[ecra=cComum] (c1) at (3.9,-8.3)  {Notificações\rota{painel no cabeçalho}};
  \node[ecra=cComum] (c2) at (7.8,-8.3)  {Definições\rota{/settings}};
  \node[ecra=cComum] (c3) at (11.7,-8.3) {Alternar tema e terminar sessão};

  % ---------- Ligações ----------
  % Login: encaminha para a página inicial de cada perfil
  \draw[ligacao] ($(p1.north)+(0.9,0)$) -- (0.9,2.1) -| ($(a1.north)+(1.0,0)$);
  \draw[ligacao] (0.9,2.1) -| ($(s1.north)+(1.0,0)$);
  \draw[ligacao] (0.9,2.1) -| ($(u1.north)+(1.0,0)$);
  \node[rotulo, anchor=south west] at (1.0,2.1) {após login, conforme o perfil};

  % Administrador
  \draw[ligacao] (a1.west) -- ++(-0.25,0) |- (a2.west);
  \draw[ligacao] (a1.west) -- ++(-0.25,0) |- (a3.west);
  \draw[ligacao] (a3) -- (a4);

  % Supervisor
  \draw[ligacao] (s1.west) -- ++(-0.25,0) |- (s2.west);
  \draw[ligacao] (s1.west) -- ++(-0.25,0) |- (s4.west);
  \draw[ligacao] (s1.west) -- ++(-0.25,0) |- (s5.west);
  \draw[ligacao] (s2) -- (s3);
  \draw[atalho] (s2.east) -- ++(0.22,0) |- (s4.east);

  % Supervisionado
  \draw[ligacao] (u1) -- (u2);
  \draw[ligacao] (u2) -- (u3);
  \draw[ligacao] (u1.west) -- ++(-0.25,0) |- (u4.west);
  \draw[ligacao] (u1.west) -- ++(-0.25,0) |- (u5.west);

  % ---------- Áreas (fundo) ----------
  \begin{scope}[on background layer]
    \node[area=gInkSoft, dashed, fill=none, fit=(la)(lu)(a1)(u5)(c1)(c3), inner sep=14pt] (aut) {};
    \node[rotulo, font=\scriptsize\bfseries, anchor=north east] at (aut.south east) {Área autenticada};
    \node[area=cPublico, fit=(lp)(p1)(p2)] {};
    \node[area=cAdmin, fit=(la)(a1)(a4)] {};
    \node[area=cSupervisor, fit=(ls)(s1)(s5)] {};
    \node[area=cAluno, fit=(lu)(u1)(u5)] {};
    \node[area=cComum, fit=(lc)(c1)(c3)] {};
  \end{scope}
\end{tikzpicture}
}
    \caption{Site map do GamP-SCPeriop. Cada caixa é um ecrã, com a rota por baixo; as setas cheias são a navegação principal e a tracejada um atalho.}
    \label{fig:sitemap}
\end{figure}

\section{Organização dos conteúdos}

Os conteúdos seguem a hierarquia do domínio: um percurso tem módulos, e um módulo tem componentes. O site map reproduz essa hierarquia nas duas áreas que a mostram:

\begin{itemize}
    \item Na área do supervisor, o painel lista os percursos; cada percurso abre a lista de inscritos, e cada inscrito abre a sua ficha de avaliação.
    \item Na área do supervisionado, \emph{Os meus percursos} lista os percursos; cada percurso abre os seus módulos, e cada módulo as suas componentes e documentos.
\end{itemize}

\section{Áreas e diferenças de acesso}

A Tabela~\ref{tab:acessos} indica, para cada área, quem lhe pode aceder e por onde se entra.

\begin{table}[H]
    \centering
    \small
    \begin{tabularx}{\linewidth}{llX}
        \toprule
        \textbf{Área} & \textbf{Acesso} & \textbf{Entrada e ecrãs principais} \\
        \midrule
        Pública & Qualquer visitante & Entrada e login (\texttt{/}); página não encontrada. \\
        \midrule
        Administrador & Perfil Administrador & Painel de administração; gestão de utilizadores; moldes de percurso e respetivo construtor. \\
        \midrule
        Supervisor & Perfil Supervisor & Painel com os percursos; gestão de cada percurso e dos inscritos; ficha de avaliação; construtor de percursos; moldes. \\
        \midrule
        Supervisionado & Perfil Supervisionado & Os meus percursos; detalhe do percurso e do módulo; conquistas; relatório de avaliação em PDF. \\
        \midrule
        Comum & Os três perfis & Notificações e alternância de tema no cabeçalho; definições (alterar a palavra-passe); terminar sessão. \\
        \bottomrule
    \end{tabularx}
    \caption{Áreas da aplicação e respetivo acesso.}
    \label{tab:acessos}
\end{table}

\section{Ligações entre ecrãs}

Além da navegação hierárquica, há três tipos de ligação que atravessam o site map:

\begin{itemize}
    \item \textbf{Construtor de percursos.} É o mesmo ecrã em três modos. O administrador usa-o em modo \emph{Molde}; o supervisor usa-o em modo \emph{Geral}, a partir do painel ou da gestão do percurso, e em modo \emph{Individual}, a partir de um inscrito.
    \item \textbf{Notificações.} Cada notificação leva diretamente ao ecrã a que diz respeito: uma nova avaliação abre o detalhe do percurso, e uma nova badge abre as conquistas.
    \item \textbf{Moldes.} O ecrã de moldes é partilhado pelo administrador, que gere os moldes-base, e pelo supervisor, que os consulta para criar percursos.
\end{itemize}

Os ecrãs abaixo da página inicial de cada perfil têm um botão para voltar ao ecrã de cima, e o cabeçalho é igual em todos os ecrãs autenticados.
